% CMcal supplied by the original eucal package already loaded below.
\usepackage[margin=1in, centering]{geometry}
\usepackage{amsthm}
\makeatletter
\let\@fnsymbol\@arabic
\makeatother
\usepackage[table]{xcolor}
\usepackage[normalem]{ulem}
%\usepackage{bbm}
\definecolor{TSUYUKUSA}{RGB}{46, 169, 223}
\definecolor{KURENAI}{RGB}{203, 27, 69}
\usepackage[pdfstartview=FitH,pdfpagemode=UseNone,colorlinks=true,citecolor=KURENAI,linkcolor=TSUYUKUSA,backref=page,linktocpage=true]{hyperref}


\newcommand{\mytitle}{TODO}

%\input{preamble.tex}
\usepackage{hyperref}
\usepackage{fullpage}
\usepackage{amssymb}
\usepackage{mathtools}
\usepackage{amsthm}
\usepackage{eucal}
\usepackage{dsfont}
\usepackage{makecell}
%\usepackage[T1]{fontenc}
%\usepackage{lmodern}
%\usepackage[stretch=10,shrink=10]{microtype}
\usepackage[capitalise]{cleveref}
%\usepackage{color,soul}
\usepackage{tikz}
\usepackage{thm-restate}
\usepackage{subfig}
\usepackage{framed}
\usepackage{environ}

\usepackage{enumitem}
\usepackage{algorithmic}
\usepackage{float}

\NewEnviron{Anonymous}{\ifx\AnonymousSwitch\undefined\BODY\fi}

%% new section font style
\makeatletter
\renewcommand{\section}{\@startsection {section}{1}{\z@}%
             {-3.5ex \@plus -1ex \@minus -.2ex}%
             {2.3ex \@plus .2ex}%
             {\normalfont\Large\scshape\bfseries}}
\renewcommand{\subsection}{\@startsection{subsection}{2}{\z@}%
             {-3.25ex\@plus -1ex \@minus -.2ex}%
             {1.5ex \@plus .2ex}%
             {\normalfont\large\scshape\bfseries}}
\renewcommand{\subsubsection}{\@startsection{subsubsection}{2}{\z@}%
             {-3.25ex\@plus -1ex \@minus -.2ex}%
             {1.5ex \@plus .2ex}%
             {\normalfont\normalsize\scshape\bfseries}}
\makeatother

\allowdisplaybreaks[1]

% Mathcal

\def\A{\CMcal{A}}
\def\B{\CMcal{B}}
\def\C{\CMcal{C}}
\def\D{\CMcal{D}}
\def\E{\CMcal{E}}
\def\F{\CMcal{F}}
\def\G{\CMcal{G}}
\def\H{\CMcal{H}}
\def\J{\CMcal{J}}
\def\K{\CMcal{K}}
\def\L{\CMcal{L}}
\def\M{\CMcal{M}}
\def\N{\CMcal{N}}
\def\O{\CMcal{O}}
\def\P{\CMcal{P}}
\def\Q{\CMcal{Q}}
\def\R{\CMcal{R}}
\def\S{\CMcal{S}}
\def\T{\CMcal{T}}
\def\U{\CMcal{U}}
\def\V{\CMcal{V}}
\def\W{\CMcal{W}}
\def\X{\CMcal{X}}
\def\Y{\CMcal{Y}}
\def\Z{\CMcal{Z}}
\def\ve{\varepsilon}
\def\xtilde{\tilde{X}}
\def\ytilde{\tilde{Y}}

% theorem, lemma, fact, etc.

\theoremstyle{plain}
\newtheorem{theorem}{Theorem}[section]
\newtheorem{lemma}[theorem]{Lemma}
\newtheorem{prop}[theorem]{Proposition}
\newtheorem{cor}[theorem]{Corollary}
\newtheorem*{maintheorem}{\Cref{thm:main}}
\theoremstyle{definition}
\newtheorem{definition}[theorem]{Definition}
\newtheorem{claim}[theorem]{Claim}
\newtheorem{problem}[theorem]{Problem}
\newtheorem{remark}[theorem]{Remark}
\newtheorem{fact}[theorem]{Fact}
\newtheorem{example}[theorem]{Example}


% basic math

\newcommand {\minusspace} {\: \! \!}
\newcommand {\smallspace} {\: \!}
\newcommand {\fn} [2] {\ensuremath{ #1 \minusspace \br{ #2 } }}
\newcommand {\Fn} [2] {\ensuremath{ #1 \minusspace \Br{ #2 } }}
\newcommand{\reals}{{\mathbb R}}
\newcommand {\set} [1] {\ensuremath{ \left\lbrace #1 \right\rbrace }}
\newcommand {\ceil} [1] {\ensuremath{ \left\lceil #1 \right\rceil }}
\newcommand {\floor} [1] {\ensuremath{ \left\lfloor #1 \right\rfloor }}
\newcommand {\innerproduct} [2] {\ensuremath{\left \langle #1 , #2 \right \rangle}}

% Fields
% \newcommand{\bF}{\mathbb{F}}

% brackets, norm, etc.

\newcommand{\normthree}[1]{{\left\vert\kern-0.25ex\left\vert\kern-0.25ex\left\vert #1 \right\vert\kern-0.25ex\right\vert\kern-0.25ex\right\vert}}
\newcommand {\br} [1] {\ensuremath{ \left( #1 \right) }}
\newcommand {\Br} [1] {\ensuremath{ \left[ #1 \right] }}
\newcommand {\cbr} [1] {\ensuremath{ \left\lbrace #1 \right\rbrace }}
%\newcommand {\norm} [1] {\ensuremath{ \left\| #1 \right\| }}
\DeclarePairedDelimiter{\norm}{\lVert}{\rVert}
\newcommand {\normsub} [2] {\ensuremath{ \norm{#1}_{#2} }}
\newcommand {\onenorm} [1] {\normsub{#1}{1}}
\newcommand {\twonorm} [1] {\normsub{#1}{2}}
\newcommand {\ab} [1] {\ensuremath{ \left| #1 \right| }}
\DeclarePairedDelimiter{\abs}{\lvert}{\rvert}

\newcommand {\fidelity} [2] {\mathrm{F}\br{#1,#2}}
\newcommand {\bra} [1] {\ensuremath{ \left\langle #1 \right| }}
\newcommand {\ket} [1] {\ensuremath{ \left| #1 \right\rangle }}
\newcommand {\ketbratwo} [2] {\ensuremath{ \left| #1 \middle\rangle \middle\langle #2 \right| }}
\newcommand {\ketbra} [1] {\ketbratwo{#1}{#1}}
\newcommand {\braket} [2] {\ensuremath{\left \langle #1 | #2 \right \rangle}}


% probability and information theory



%\newcommand {\defeq} {\ensuremath{ \stackrel{\mathrm{def}}{=} }}
\newcommand {\defeq} {\ensuremath{ = }}
\newcommand {\mutinf} [2] {\fn{\mathrm{I}}{#1 \smallspace : \smallspace #2}}
\newcommand {\condmutinf} [3] {\mutinf{#1}{#2 \smallspace \middle\vert \smallspace #3}}
\newcommand {\prob} [1] {\Fn{\Pr\,}{#1}}
\newcommand {\relent} [2] {\fn{\mathbb{D}}{#1 \middle\| #2}}
\newcommand {\rminent} [2] {\fn{\mathrm{D}_{\infty}}{#1 \middle\| #2}}
\DeclareMathOperator*{\bigE}{\mathbb{E}}
\newcommand {\expec} [2] {\Fn{\bigE_{\substack{#1}}}{#2}}
\newcommand{\nnorm}[1]{{\left\vert\kern-0.25ex\left\vert\kern-0.25ex\left\vert #1
		\right\vert\kern-0.25ex\right\vert\kern-0.25ex\right\vert}}

% matrix theory

\newcommand {\Tr} {\ensuremath{ \mathrm{Tr} }}
\newcommand {\partrace} [2] {\fn{\Tr_{#1}}{#2}}
\newcommand {\adjoint} [1] {\ensuremath{#1^{*}}}
\newcommand {\id} {\ensuremath{\mathds{1}}}
\newcommand {\unitary} [1] {\ensuremath{ \mathbf{#1} }}
\newcommand{\mapvec}[1]{\text{vec}\br{{#1}}}
\newcommand{\conjugate}[1]{\overline{#1}}



% communication complexity
\newcommand{\qcc}[2]{\mathrm{Q}^{\text{ent},\text{A}\rightarrow\text{B}}_{#1}\br{#2}}
\newcommand{\dqcc}[3]{\mathrm{Q}^{\text{ent},\text{A}\rightarrow\text{B},#1}_{#2}\br{#3}}
\newcommand {\Prot}{\mathcal{P}}


%miscellaneous

\newcommand {\email} [1] {\href{mailto:#1}{\texttt{#1}}}
\newcommand {\ie} {i.e.,}
\newcommand {\myqedhere} {\tag*{\qedhere}}
\newcommand{\rev}[1]{\textcolor{red}{#1}}

\newcommand {\psd}{\mathcal{D}}


\newcommand {\authorblock} [3] {
	\begin{minipage}[t]{0.3\linewidth}
		\centering
		{#1}\\[0.8ex]
		{\footnotesize {#2}\\[-0.7ex]
		\email{#3}}
	\end{minipage}\vspace{1ex}
}

%\hypersetup{
%	pdfstartview={FitH},
%	pdfdisplaydoctitle={true},
%	breaklinks={true},
%	bookmarksopen={true},
%	bookmarksnumbered={false},
%	pdftitle={\mytitle}
%%	pdfauthor={\Anurag, \Rahul, \Priyanka, \Ala, \Penghui}
%}


% tikz
\usetikzlibrary{calc}
\tikzset{meter/.append style={draw, inner sep=10, rectangle, font=\vphantom{A}, minimum width=30, line width=.8,
 path picture={\draw[black] ([shift={(.1,.3)}]path picture bounding box.south west) to[bend left=50] ([shift={(-.1,.3)}]path picture bounding box.south east);\draw[black,-latex] ([shift={(0,.1)}]path picture bounding box.south) -- ([shift={(.3,-.1)}]path picture bounding box.north);}}}

\newenvironment{Algorithm}{\begin{framed}}{\end{framed}}

% Cref
\crefname{prop}{Proposition}{Propositions}
  
%\input{commands.tex}
% qml add

\newcommand{\inv}[1]{\ensuremath{#1^{-1}}}
\newcommand{\pos}[1]{\ensuremath{#1^{+}}}
\newcommand{\sqrtinv}[1]{\ensuremath{{#1}^{-\frac{1}{2}}}}
\renewcommand{\vec}[1]{\stackrel{\longrightarrow}{#1}}
\newcommand{\posint}{{\mathbb Z}_{>0}}
\newcommand{\e}{\mathrm{e}}
\renewcommand{\d}{\mathrm{d}}
\renewcommand{\i}{\mathrm{i}}
%\newcommand{\reals}{{\mathbb R}}
\newcommand{\complex}{{\mathbb C}}
\newcommand{\randP}{\ensuremath{\mathbf{P}}}
\newcommand{\randb}{\ensuremath{\mathbf{b}}}
\newcommand{\randg}{\ensuremath{\mathbf{g}}}
\newcommand{\randh}{\ensuremath{\mathbf{h}}}
\newcommand{\randz}{\ensuremath{\mathbf{z}}}
\newcommand{\randx}{\ensuremath{\mathbf{x}}}
\newcommand{\randy}{\ensuremath{\mathbf{y}}}
\newcommand{\randu}{\ensuremath{\mathbf{u}}}
\newcommand{\randv}{\ensuremath{\mathbf{v}}}
\newcommand{\randw}{\ensuremath{\mathbf{w}}}
\newcommand{\randf}{\ensuremath{\mathbf{f}}}
\newcommand{\randA}{\ensuremath{\mathbf{A}}}
\newcommand{\randB}{\ensuremath{\mathbf{B}}}
\newcommand{\randC}{\ensuremath{\mathbf{C}}}
\newcommand{\randD}{\ensuremath{\mathbf{D}}}
\newcommand{\randM}{\ensuremath{\mathbf{M}}}
\newcommand{\randX}{\ensuremath{\mathbf{X}}}
\newcommand{\randa}{\ensuremath{\mathbf{a}}}
\newcommand{\randPi}{\ensuremath{\mathbf{\Pi}}}
\newcommand{\herspace}[1]{\ensuremath{\H_m^{\otimes #1}}}
\newcommand{\randspace}[2]{L^2\br{\herspace{#1},\gamma_{#2}}}
\newcommand{\infi}[1]{\ensuremath{\mathrm{Inf}_i}\br{#1}}
\newcommand{\varinfi}[1]{\ensuremath{\mathrm{VarInf}_i}\br{#1}}
\newcommand{\inffj}[1]{\ensuremath{\mathrm{VarInf}_{f,j}}\br{#1}}
\newcommand{\infrfj}[1]{\ensuremath{\mathrm{VarInf}_{\randf,j}}\br{#1}}
\newcommand{\invf}[1]{\ensuremath{f^{-1}\br{#1}}}
\newcommand{\bits}[1]{\ensuremath{\set{-1,1}^{#1}}}
\newcommand{\usim}{\ensuremath{\sim_{\mathrm{u}}}}
\newcommand {\var} [1] {\mathrm{Var}\Br{#1}}
\newcommand{\gval}[3]{\ensuremath{\mathrm{val}_{#1}\br{\set{#2},\set{#3}}}}
\newcommand{\fourtuples}{\ensuremath{\br{x,y,a,b}\in\calX\times\calY\times\calA\times\calB}}
\newcommand{\sigmarange}[1]{\ensuremath{\sigma\in[m^2]_{\geq0}^{#1}}}
\newtheorem*{theorem*}{Theorem}

%%%% dyj
\newcommand{\TM} {\mathcal{M}}
\newcommand{\poly} {\ensuremath{\operatorname{poly}}}
\newcommand{\polylog} {\ensuremath{\operatorname{polylog}}}
\newcommand{\polyloglog} {\ensuremath{\operatorname{polyloglog}}}
\newcommand{\ep} {\ensuremath{\varepsilon}}
\newcommand{\resp} {resp.}

\newcommand{\ccsm} {\ensuremath{C_{sm}}}
\newcommand{\dmcc} {\ensuremath{C_{pt}}}
\newcommand{\dlbo} {\ensuremath{{D}}}
\newcommand{\dldd} {\ensuremath{{d}}}
\newcommand{\wireal} {\ensuremath{{s_{w}}}}
\newcommand{\epsm} {\ensuremath{\ep^\prime}}
\newcommand{\epreg} {\ensuremath{{\delta}}}
\newcommand{\epdr} {\ensuremath{\delta}}
\newcommand{\epfn} {\ensuremath{\delta_{fn}}}
\newcommand{\eprd} {\ensuremath{\ep_{rd}}}
\newcommand{\redt}[1] {\textcolor{red}{#1}} % make inside tex red
\newcommand{\pseudostrategy} {pseudo-strategy}
\newcommand{\psdtesttime} {\ensuremath{\operatorname{PSDTIME}}}
\newcommand{\nonlocalgametime} {\ensuremath{\operatorname{NLGTIME}}}
\newcommand{\Pss} {\ensuremath{P^{\,x}_a}}
\newcommand{\Qss} {\ensuremath{Q^{\,y}_b}}
\newcommand{\Pwss} {\ensuremath{\tilde{P}^{\,x}_a}}
\newcommand{\Qwss} {\ensuremath{\tilde{Q}^{\,y}_b}}
\newcommand{\Pnn}[1] {\ensuremath{{P^{\,x, (#1)}_a}}}
\newcommand{\Qnn}[1] {\ensuremath{{Q^{\,y, (#1)}_b}}}
\newcommand{\Pff} {\ensuremath{\widehat{P}^{\,x}_a(\sigma)}}
\newcommand{\Qff} {\ensuremath{\widehat{Q}^{\,y}_b(\sigma)}}
\newcommand{\Pfn}[1] {\ensuremath{{{\widehat{P}}^{\,x, {(#1)}}_a}(\sigma)}}
\newcommand{\Qfn}[1] {\ensuremath{{{\widehat{Q}}^{\,y, {(#1)}}_b}(\sigma)}}

%% Answer reduction
\newcommand{\x}{\otimes}
\newcommand{\xp}[1]{^{\otimes #1}}
\newcommand{\op}{\oplus}

\newcommand{\ct}{^{\dagger}}
\newcommand{\tp}{^\intercal}

\newcommand{\calH}{\mathcal{H}}
\newcommand{\calX}{\mathcal{X}}
\newcommand{\calY}{\mathcal{Y}}
\newcommand{\calA}{\mathcal{A}}
\newcommand{\calB}{\mathcal{B}}
\newcommand{\calU}{\mathcal{U}}
\newcommand{\calL}{\mathcal{L}}
\newcommand{\calF}{\mathcal{F}}
\newcommand{\calG}{\mathcal{G}}
\newcommand{\calR}{\mathcal{R}}
\newcommand{\whA}{\widehat{A}}
\newcommand{\whB}{\widehat{B}}
\newcommand{\whM}{\widehat{M}}
\newcommand{\whN}{\widehat{N}}

\newcommand{\appd}[1]{\approx_{#1}}

\newcommand{\1}{\mathbb{1}} % identity matrix

\newcommand{\fC}{\mathfrak{C}}

\newcommand{\bF}{\mathbb{F}}
\newcommand{\bZ}{\mathbb{Z}}

\newcommand{\ba}{\pmb{a}}
\newcommand{\bb}{\pmb{b}}
\newcommand{\bc}{\pmb{c}}
\newcommand{\bd}{\pmb{d}}
\newcommand{\be}{\pmb{e}}
\newcommand{\bw}{\pmb{w}}
\newcommand{\bx}{\pmb{x}}
\newcommand{\by}{\pmb{y}}
\newcommand{\bz}{\pmb{z}}
\newcommand{\bI}{\pmb{I}}
\newcommand{\bJ}{\pmb{J}}
\newcommand{\bT}{\pmb{T}}
\newcommand{\bi}{\pmb{i}}
\newcommand{\bj}{\pmb{j}}
\newcommand{\bU}{\pmb{U}}
\newcommand{\pmbF}{\pmb{F}}
\newcommand{\bS}{\mathbbmtt{S}}

\newcommand{\tO}{\tilde{O}}
\newcommand{\tP}{\tilde{P}}
\newcommand{\tQ}{\tilde{Q}}
\newcommand{\tM}{\tilde{M}}
\newcommand{\tN}{\tilde{N}}
\newcommand{\tA}{\tilde{A}}
\newcommand{\tB}{\tilde{B}}
\newcommand{\tG}{\tilde{G}}
\newcommand{\tH}{\tilde{H}}
\newcommand{\tX}{\tilde{X}}
\newcommand{\tZ}{\tilde{Z}}
\newcommand{\tU}{\tilde{U}}
\newcommand{\tW}{\tilde{W}}
\newcommand{\tx}{\tilde{x}}
\newcommand{\ta}{\tilde{a}}
\newcommand{\tpsi}{\tilde{\psi}}

%\DeclarePairedDelimiter{\set}{\lbrace}{\rbrace}
\newcommand{\se}{\sqrt{\epsilon}}
\newcommand{\qe}{\epsilon^{1/4}}
\newcommand{\PCPP}{\mathrm{PCPP}}
\newcommand{\Unif}{\mathsf{Unif}}
\newcommand{\ipt}{\mathsf{input}}
\newcommand{\Enc}{\mathsf{Enc}}
\newcommand{\Dec}{\mathsf{Dec}}
\newcommand{\Alg}{\mathsf{Alg}}
\newcommand{\MIP}{\mathrm{MIP}}
\newcommand{\RE}{\mathrm{RE}}
\newcommand{\val}{\mathrm{val}}
\newcommand{\NEXP}{\mathrm{NEXP}}
\newcommand{\NP}{\mathrm{NP}}


\DeclareMathOperator*{\argmax}{argmax}
\DeclareMathOperator*{\argmin}{arg\,min}

% Comments
\newcommand{\honghao}[1]{\textcolor{orange}{Honghao: {#1}}}

\newcommand{\anand}[1]{\textcolor{red}{Anand: {#1}}}

\newcommand{\mlqin}[1]{\textcolor{blue}{mlqin: {#1}}}


\setlength\parindent{20pt}
\allowdisplaybreaks
